\section{Hurwitz zeta, negative polygamma, and the analytic hyperfactorial}\label{sec:hurwitz}

\subsection{A constructive Hurwitz-zeta derivative}

For \(\Re z>0\), the Hurwitz zeta function has the fixed-order
Euler--Maclaurin representation
\begin{align}
 \zeta(s,z)={}&\sum_{k=0}^{N-1}(z+k)^{-s}
 +\frac{(z+N)^{1-s}}{s-1}+\frac12(z+N)^{-s}\notag\\
 &+\sum_{r=1}^{M}
 \frac{B_{2r}}{(2r)!}(s)_{2r-1}(z+N)^{-s-2r+1}
 +R_{M,N}(s,z),
 \label{eq:hurwitz-EM}
\end{align}
where \((s)_j=s(s+1)\cdots(s+j-1)\).  On a compact box in \(s,z\), with
\(\Re(z+N)\) positively bounded, the usual periodic-Bernoulli integral gives
an explicit rational majorant for \(R_{M,N}\) and its derivative in \(s\).
For fixed \(M\), all coefficients are rational; powers and logarithms use the
branch path supplied with \(z\).  Thus \(\zeta'(-1,z)\) has a
primitive-recursive fast name on every compact subset of \(\Re z>0\).

The fundamental finite shift relation is
\begin{equation}\label{eq:hurwitz-shift}
 \zeta(s,z+1)=\zeta(s,z)-z^{-s}.
\end{equation}
Differentiating at \(s=-1\) gives
\begin{equation}\label{eq:hurwitz-derivative-shift}
 \zeta'(-1,z+1)-\zeta'(-1,z)=z\Log z.
\end{equation}
The differentiation can be avoided formally: apply a symmetric rational
difference quotient in \(s\) to the exact shift identity, bound the Taylor
remainder uniformly, and then let the precision parameter increase.  This is a
schema of finite rational statements.

\subsection{Vardi's formula}

Define on \(\Re z>0\)
\begin{equation}\label{eq:vardi-rhs}
 V(z)=\frac1{12}-\zeta'(-1,z)+(z-1)\Log\Gamma(z)-\log A.
\end{equation}
Using \cref{eq:hurwitz-derivative-shift},
\(\Log\Gamma(z+1)=\Log z+\Log\Gamma(z)\), and
\(\log A=1/12-\zeta'(-1)\), one finds
\begin{align*}
 V(z+1)-V(z)
 &=-z\Log z+z\Log\Gamma(z+1)-(z-1)\Log\Gamma(z)\\
 &=\Log\Gamma(z).
\end{align*}
Also \(V(1)=0\).  The difference
\(V(z)-\Log G(z)\) is therefore one-periodic.  Euler--Maclaurin applied to
\cref{eq:vardi-rhs}, or the asymptotic theorem below, shows that this difference
tends to zero as \(\Re z\to+\infty\) in a fixed horizontal strip.  Periodicity
then forces it to be identically zero.

\begin{theorem}[Hurwitz-zeta derivative formula]\label{thm:vardi}
For \(\Re z>0\), with logarithms continued from the positive real axis,
\begin{equation}\label{eq:vardi}
 G(z)=\frac{
 \exp(1/12-\zeta'(-1,z))\,\Gamma(z)^{z-1}}{A}.
\end{equation}
Equivalently,
\[
 \Log G(z)=\frac1{12}-\zeta'(-1,z)
 +(z-1)\Log\Gamma(z)-\log A.
\]
\end{theorem}

\begin{proof}[Effective periodicity argument]
Given \(z\) in a certified compact subset of \(\Re z>0\) and precision \(p\),
choose an integer \(m\) so that the explicit Euler--Maclaurin remainders for
\(V(z+m)\) and \(\Log G(z+m)\) sum to less than \(2^{-p}\).  Repeated use of
the exact shift equations gives
\[
 V(z)-\Log G(z)=V(z+m)-\Log G(z+m).
\]
The right side lies in the dyadic ball of radius \(2^{-p}\) around zero.  This
is the precision instance of equality.  The shift count and every remainder
bound are primitive recursive in the query data.
\end{proof}

This proves MathWorld equation (4), including its domain qualification.

\subsection{Glaisher's constant as an iPRA name}

Let
\[
 \hyp(N)=\prod_{k=1}^{N}k^k.
\]
Euler--Maclaurin for \(\sum_{k=1}^N k\log k\) gives
\begin{equation}\label{eq:hyper-asymp}
 \hyp(N)=A\,
 N^{N^2/2+N/2+1/12}\e^{-N^2/4}
 \bigl(1+\varepsilon_N\bigr),
 \qquad |\varepsilon_N|\le \frac{C}{N^2}
\end{equation}
with an explicit rational \(C\) once a fixed truncation order is selected.
Thus
\begin{equation}\label{eq:A-limit}
 \log A=\lim_{N\to\infty}
 \left[
 \sum_{k=1}^{N}k\log k
 -\left(\frac{N^2}{2}+\frac N2+\frac1{12}\right)\log N
 +\frac{N^2}{4}
 \right]
\end{equation}
has a primitive-recursive modulus.  Evaluating the same Euler--Maclaurin
constant through the analytic continuation of zeta gives
\begin{equation}\label{eq:A-zeta-again}
 \log A=\frac1{12}-\zeta'(-1).
\end{equation}

For formalization one may choose \cref{eq:A-limit} as the definition of the
name \(A\), because it uses only finite sums, positive logarithms, and an
explicit remainder.  Equation \cref{eq:A-zeta-again} then becomes a theorem.
Alternatively, if the zeta module is already available, the right side of
\cref{eq:A-zeta-again} is shorter.  Both names are proved equal by comparing
their finite Euler--Maclaurin approximants.

\subsection{The analytic hyperfactorial or K-function}

To avoid confusing the K-function with the multiplication constant \(K(n)\),
write \(\mathcal H(z)\) for the analytic continuation of the hyperfactorial,
normalized by
\begin{equation}\label{eq:analytic-hyper-recurrence}
 \mathcal H(1)=1,
 \qquad
 \mathcal H(z+1)=z^z\mathcal H(z).
\end{equation}
On a branch domain define
\begin{equation}\label{eq:analytic-hyper-def}
 \mathcal H(z)=
 \frac{\exp((z-1)\Log\Gamma(z))}{G(z)}.
\end{equation}
Then \(\mathcal H(1)=1\), and
\begin{align*}
 \frac{\mathcal H(z+1)}{\mathcal H(z)}
 &=\frac{\exp(z\Log\Gamma(z+1))}{G(z+1)}
   \frac{G(z)}{\exp((z-1)\Log\Gamma(z))}\\
 &=\exp(z\Log z)=z^z.
\end{align*}
Therefore \(\mathcal H\) is the normalized analytic K-function and
\begin{equation}\label{eq:analytic-hyper-relation}
 \mathcal H(z)G(z)=\exp((z-1)\Log\Gamma(z)).
\end{equation}
At \(z=n\in\N_{>0}\), the recurrence gives
\(\mathcal H(n)=\hyp(n-1)\), and \cref{eq:analytic-hyper-relation} reduces to
the exact identity \cref{eq:hyper-cross}.

The formal primary statement is the exponentiated, cross-multiplied equality.
A global claim ``for all complex \(z\)'' must be read on the logarithmic
cover, because \(z^z\) and \(\Gamma(z)^{z-1}\) are branch dependent.  On any
fixed path domain, the continuation algorithm of \cref{sec:ipra} supplies the
needed name.

\subsection{Negative-order polygamma form}

A normalized negative-order polygamma \(\psi_{-2}\) can be characterized by
\begin{equation}\label{eq:negpoly-rec}
 \psi_{-2}(z+1)-\psi_{-2}(z)
 =z\Log z-z+\frac12\log(2\pi),
 \qquad
 \psi_{-2}(1)=\frac12\log(2\pi),
\end{equation}
plus the standard eventual higher-convexity or asymptotic normalization.  Set
\begin{equation}\label{eq:negpoly-expression}
 W(z)=-\binom z2+(z-1)\Log\Gamma(z)
      +\frac z2\log(2\pi)-\Log G(z).
\end{equation}
Using the Gamma and Barnes recurrences, a finite calculation shows that \(W\)
satisfies \cref{eq:negpoly-rec}; the asymptotic module supplies the normalization.
Hence \(W=\psi_{-2}\), and
\begin{equation}\label{eq:negative-polygamma}
 \Log G(z)=
 -\binom z2+(z-1)\Log\Gamma(z)
 +\frac z2\log(2\pi)-\psi_{-2}(z).
\end{equation}
Exponentiating gives the MathWorld form
\begin{equation}\label{eq:negative-polygamma-exp}
 G(z)=(2\pi)^{z/2}
 \exp\!\left((z-1)(\Log\Gamma(z)-z/2)-\psi_{-2}(z)\right).
\end{equation}

The qualification ``normalized'' is essential.  An indefinite second
antiderivative is determined only up to an affine function; changing that
normalization changes the right side.  The finite recurrence and value in
\cref{eq:negpoly-rec} remove the ambiguity and provide the object-iPRA target.

\begin{roadmap}[Lean modules \code{HurwitzZeta/DerivativeMinusOne} and
\code{BarnesG/Vardi}]
Implement \cref{eq:hurwitz-EM} for a fixed external Euler--Maclaurin order, with
both value and \(s\)-derivative remainder checkers.  Prove the shift relation
first at finite sums, then the derivative shift by a certified difference
quotient.  The Vardi theorem is a generic lemma: two named solutions of the
same first-order difference equation whose difference has an effective limit
zero on right shifts are equal.  Reuse that lemma for the negative-polygamma
and hyperfactorial normalizations.
\end{roadmap}
